% Part: counterfactuals
% Chapter: minimal-change-semantics
% Section: antecedent-strengthening

\documentclass[../../../include/open-logic-section]{subfiles}

\begin{document}

\olfileid{cnt}{min}{agg}

\olsection{Penguwatan Anteseden}

``Penguwatan anteseden'' ngrujuk marang inferensi
$!A \lif !C \Entails (!A \land !B) \lif !C$. Inferensi iki sah
kanggo kondisional material, nanging ora sah kanggo kontrafaktual.
Upamane bener yen saupama aku nggesek penthol korek iki, penthol
iku mesthi bakal murub. (Tegese, ora ana kang salah ing penthol
korek utawa lumahing bungkus korek, aku ora bakal mecah penthol
korek, lan sapiturute.) Nanging ora bener yen saupama aku
nggesek penthol korek iki ing angkasa njaba, penthol iku mesthi
bakal murub. Mula inferensi iki ora sah:
\begin{quote}
  Saupama penthol korek digesek, penthol iku mesthi bakal murub.

  Mula, saupama penthol korek digesek ing angkasa njaba,
  penthol iku mesthi bakal murub.
\end{quote}

Panjelasan Lewis--Stalnaker tumrap kondisional nerangake iki:
donya paling cedhak kang ing kono aku nggesek penthol korek lan
nindakake ing angkasa njaba luwih adoh saka donya aktual tinimbang
donya paling cedhak kang ing kono aku nggesek penthol korek.
Mula, sanajan penthol korek murub ing donya kapindho, ora murub
ing donya kapisan. Iki cocog karo kang dikarepake.

\begin{ex}
  Semantik bola nuduhake yen inferensi iku ora sah, yaiku $p
  \cif r \Entails/ (p \land q) \cif r$. Gatekna model
  $\mModel{M} = \tuple{W, O, V}$ kanthi $W = \{w, w_1, w_2\}$,
  $O_w = \{\{w\}, \{w, w_1\}, \{w, w_1, w_2\}\}$,
  $V(p) = \{w_1, w_2\}$, $V(q) = \{w_2\}$, lan $V(r) = \{w_1\}$.
  Ana bola kang ngakoni $p$, yaiku $S = \{w, w_1\}$, lan
  $p \lif r$ bener ing kabeh donya ing kono, mula
  $\mSat{M}{p \cif r}[w]$. Bola siji-sijine kang ngakoni
  $(p \land q)$ yaiku $S' = \{w, w_1, w_2\}$, nanging
  $\mSat/{M}{(p \land q) \lif r}[w_2]$, mula
  $\mSat/{M}{(p \land q) \cif r}[w]$
  (delengen \olref{fig:antecedent-strengthening}).

\begin{figure}
\begin{center}
\begin{tikzpicture}[layerwidth=1.5,scale=.6]\tiny
  \spheresystem{3}
  \propositionintersect{-10}{3}{30}{5}
  \begin{scope}
    \clip plot[smooth,tension=1] coordinates
          {(0,4.5) (30:.95) (4.5,-3.5)} -- (4.5,4.5) ;
    \spherefill{2}
  \end{scope}
  \draw[proposition] plot[dashed,smooth,tension=1] coordinates
          {(0,4.5) (30:.95) (4.5,-3.5)} ;
  \proposition{30}{2}{30}{5}
  \path (0,0) node[world] {$w$};
  \spherepos{30}{2}{node[world] {$w_1$}}
  \spherepos{-10}{3}{node[world] {$w_2$}}
  \spherepos{-10}{4}{node {$q$}}
  \spherepos{30}{4}{node {$r$}}
  \draw (1,4.5) node[below] {$p$};
\end{tikzpicture}
\caption{Tuladha panyanggah tumrap penguwatan anteseden}
\ollabel{fig:antecedent-strengthening}
\end{center}
\end{figure}
\end{ex}

\end{document}
